\section{The prescribed lex-plus-powers ideal}\label{sec:assembly}

We now combine the comparisons, keeping the same ideal at every
homological and internal degree.

\begin{proof}[Proof of Theorem~\ref{thm:main}]
If \(I=S\), then \(L_d=S_d\) is allowed in every degree and \(J=S\);
all quotient Betti numbers are zero. Assume \(I\ne S\).
Proposition~\ref{mono:bound} gives a monomial ideal \(M\supseteq P\)
with the Hilbert function of \(I\) and
\[
 \beta^S_{p,j}(S/I)\le\beta^S_{p,j}(S/M)
 \qquad(p\ge0,\ j\in\Z).
\]
Proposition~\ref{stab:main} gives a box-stable monomial ideal
\(H\supseteq P\), with the same Hilbert function, and
\[
 \beta^S_{p,j}(S/M)\le\beta^S_{p,j}(S/H).
\]
By Proposition~\ref{box:ideal}, the degreewise box-lex replacement
\(G=\Lambda H\) is an ideal, with the same Hilbert function.
Theorem~\ref{tor:stable} gives
\[
 \beta^S_{p,j}(S/H)\le\beta^S_{p,j}(S/G).
\]

It remains to check that \(G\) is the ideal prescribed in
\eqref{eq:prescription}.
In degree \(d\), list all monomials in decreasing lex order and remove
those belonging to \(P_d\). The remaining list is exactly
\(\mathcal B_n(d)\) in decreasing lex order.
The monomials of \(G_d\) outside \(P_d\) are the first
\[
 q_d=\dim_\C I_d-\dim_\C P_d
\]
positions of this remaining list. This integer lies between zero and
\(|\mathcal B_n(d)|\), because \(H\supseteq P\) has the same Hilbert
function as \(I\).

Every box prefix is obtained by intersecting some full-list prefix with
the box. Thus a lex-segment subspace \(L_d\) with
\(\dim_\C(L_d+P_d)=\dim_\C I_d\) exists. Enlarging that full-list prefix
to the largest one of the same resulting dimension adds only monomials
already in \(P_d\), and does not change the sum. Consequently
\(L_d+P_d=G_d\) for every \(d\), including the empty and full box cases.
Hence \(J=G\). The displayed comparisons prove all the asserted Betti
bounds for this single \(J\).
\end{proof}

The argument proves idealness and all syzygy bounds together.
The use of \(\C\) enters the commuting-form construction and supplies the
roots of unity for stabilization; no conclusion in positive
characteristic is asserted for an arbitrary regular sequence.
